%=======================================================
% ESTATUTOS DE CervanTeX,
%
% APROBADO EN LA ASAMBLEA GENERAL DE CervanTeX
% EL 14 DE SEPTIEMBRE DE 1999 EN MADRID
%
% Inscrita en el registro de asociaciones con fecha
% 14-diciembre-1999 y n\'umero de registro 165745.
%
% Ley de asociaciones: Ley 191/1964 y decreto 1440/1965.
%
% Coordinaci\'on de la redacci\'on de este documento:
% - Javier Bezos, jbezos@arrakis.es
% - Jos\'e Ra Portillo, josera@cica.es
% - Juan Luis Varona, jvarona@dmc.unirioja.es
% Escrito en plain-TeX (y con simplicidad extrema).
%=======================================================

\magnification\magstep1
\frenchspacing
\parindent=0pt % Sin sangr\'{\i}a
\parskip=4pt plus 2pt minus1pt % Separaci\'on entre p\'arrafos

% Para minimizar cortes de palabras. As\'{\i} se evita que corte mal.
\pretolerance=2000
\tolerance=2500
\hyphenpenalty=2000
\hbadness=2500 % Para evitar posibles errores de undefull al no cortar palabras

% T\'ITULOS:
\def\titulo#1#2{\par\goodbreak\medskip%
    \noindent{\bf T\'ITULO #1: #2.}%
    \par\nobreak\smallskip}

% Cap\'{\i}tulos:
\def\capitulo#1#2{\par\goodbreak\medskip%
    \noindent{\bf Cap\'{\i}tulo #1. #2.}%
    \par\nobreak\smallskip}

% Art\'{\i}culos:
\def\articulo #1.{\par\smallskip{\bf Art\'{\i}culo #1.} \thinspace}

%%%%%%%%%%%%%%%%%%% Macros para usar << y >>
\catcode`\>=\active
\catcode`\<=\active
%\font\ly=lasy10
%\def<<{\nobreak\leavevmode
%       \hbox{\ly\char40 \kern-0.19em\char40 }}
%\def>>{\nobreak\leavevmode
%       \hbox{\ly\char41 \kern-0.19em\char41 }}
% Da unas comillas iguales a las de spanishb (de donde est\'a adaptado).
% Esta definici\'on es completamente ad hoc para este archivo.
% Una alternativa para no usar la letra lasy:
\def<<{\nobreak\leavevmode
       \setbox0\hbox{$\scriptscriptstyle\langle\mkern-4mu\langle$}%
       \raise.74\dp0\box0 }
\def>>{\nobreak\leavevmode
       \setbox0\hbox{$\scriptscriptstyle\rangle\mkern-4mu\rangle$}%
       \raise.74\dp0\box0 }
%%%%%%%%%%%%%%%%%%% Fin macros para usar << y >>

\centerline{\bf ESTATUTOS DE}
\centerline{\bf Cervan\TeX{} --- Grupo de Usuarios de \TeX{} Hispanohablantes}

\medskip

\centerline{\bf 14 de septiembre de 1999}

\bigskip

%=======================================================
\titulo{I}{Aspectos generales}
%=======================================================

\articulo 1. Nombre.\par\nobreak

1. La asociaci\'on se denomina <<Cervan\TeX{} --- Grupo de Usuarios de
\TeX{} Hispano\-hablan\-tes>>. Su denominaci\'on habitual, por brevedad,
ser\'a <<Cervan\TeX>>.
Cuando no sea posible escribir <<Cervan\TeX>>, deber\'a
escribirse <<CervanTeX>>.

2. La asociaci\'on es de inter\'es p\'ublico y sin \'animo de lucro.

%------------------

\articulo 2. Domicilio.\par\nobreak

1. El domicilio social actual de Cervan\TeX{} es el Departamento de
Matem\'atica Aplicada~I de la Universidad de Sevilla, sito en la calle
Reina Victoria, n\'u\-me\-ro~2, de Sevilla.

2. El domicilio social podr\'a ser cambiado por decisi\'on
de la Junta Directiva, debiendo ser refrendado en la siguiente
Asamblea General.

%------------------

\articulo 3. \'Ambito de actuaci\'on.\par\nobreak

1. El \'ambito de actuaci\'on de Cervan\TeX{} ser\'a el territorio espa\~nol.

2. No obstante, Cervan\TeX{} procurar\'a mantenerse coordinado con otros
grupos de usuarios de \TeX{} constituidos en otros pa\'{\i}ses y destinados
a la promoci\'on y normalizaci\'on del uso de \TeX{} en espa\~nol.

%------------------

\articulo 4. Fines de la asociaci\'on.\par\nobreak

1. Los fines de la asociaci\'on son intercambiar experiencias sobre
el sistema inform\'atico de composici\'on tipogr\'afica denominado
\TeX, creado por D.~E.~Knuth, as\'{\i} como sobre sus aplicaciones,
y promoverlo de forma adecuada en el \'ambito hispanohablante.

2. Estos fines podr\'an desarrollarse mediante:

(a) La discusi\'on y la colaboraci\'on en el desarrollo de aplicaciones de
    \TeX.

(b) La presentaci\'on por los socios de propuestas para futuras versiones y
    componentes de \TeX.

(c) La colaboraci\'on y el intercambio con otros grupos de
    usuarios que compartan los fines de la asociaci\'on.

(d) Elaboraci\'on de normativas que permitan la estandarizaci\'on
    del uso de \TeX{} en espa\~nol.

(e) La promoci\'on y estimulaci\'on de publicaciones relacionadas con \TeX.

(f) La promoci\'on del uso de otros programas de distribuci\'on libre que
    puedan estar relacionados con \TeX.

(g) La representaci\'on de los intereses de los usuarios hispanohablantes
    en el \TeX{} User Group.

(h) La organizaci\'on de cursillos.

(i) La organizaci\'on de congresos.

(j) La publicidad y la emisi\'on de comunicados.

(k) El acceso de los socios a publicaciones y programas de \TeX{}
    mediante acuerdos con sus autores y propietarios legales.

(l) La formalizaci\'on de convenios de colaboraci\'on e intercambio para
    impulsar el uso de \TeX{} en todo tipo de sociedades,
    instituciones y editoriales cient\'{\i}ficas.

3. Relaciones especiales con otras asociaciones.

(a) Cervan\TeX{} mantendr\'a lazos de colaboraci\'on con otros grupos de
    usuarios de \TeX, especialmente con el \TeX{} Users Group (TUG)
    internacional, as\'{\i} como con otras asociaciones dedicadas a la
    difusi\'on y promoci\'on de {\it software\/} libre.

(b) El grupo de usuarios impulsar\'a la colaboraci\'on e intercambio,
    mediante convenios u otros mecanismos que se regulen, con asociaciones
    e instituciones dedicadas al estudio y promoci\'on del lenguaje espa\~nol.

(c) La asociacion procurar\'a mantener especiales lazos de colaboraci\'on
    con otros grupos de usuarios de \TeX{} constituidos en los \'ambitos
    geogr\'aficos donde el idioma espa\~nol sea oficial o de uso com\'un.
    Entre ellos, y especialmente, con los mencionados en el p\'arrafo 2 del
    art\'{\i}culo 3 de estos Estatutos.

%------------------

\articulo 5. Grupos de trabajo.\par\nobreak

1. El trabajo en Cervan\TeX{} se organiza en grupos de trabajo. Cada
uno de ellos funcionar\'a de forma aut\'onoma y establecer\'a sus
propios m\'etodos de trabajo.

2. Cada grupo tendr\'a un coordinador, elegido por y de entre sus
miembros.

3. Todo socio puede adscribirse libremente a los grupos de
trabajo, para lo cual deber\'a contactar con sus
respectivos coordinadores.

4. Las sugerencias para nuevos grupos deber\'an ser dirigidas a la Junta
Directiva, que evaluar\'a y, en su caso, aprobar\'a la formaci\'on del nuevo
grupo.

5. Uno de los grupos de trabajo ser\'a el encargado de elaborar las
publicaciones de Cervan\TeX. Su coordinador tendr\'a el cargo de
Editor General.

6. Ni la asociaci\'on Cervan\TeX{}, ni ninguno de sus socios, podr\'an
comercializar los programas y paquetes inform\'aticos surgidos de la
actividad de los grupos de trabajo.
En su caso, Cervan\TeX{} podr\'a, tras acuerdo de la Junta Directiva,
distribuir libremente los programas y paquetes generados en la actividad
de los grupos de trabajo, sin perjuicio del reconocimiento de la propiedad
intelectual de sus autores.

%------------------

\articulo 6. Cargos de la asociaci\'on.\par\nobreak

1. Los cargos de la asociaci\'on no recibir\'an gratificaci\'on econ\'omica
alguna.

2. En caso de que el trabajo de los cargos supere lo razonable, la
Asamblea General podr\'a acordar la contrataci\'on del personal auxiliar
que se considerase necesario. Dicho personal auxiliar estar\'a bajo las
\'ordenes del Secretario.

\bigskip

%=======================================================
\titulo{II}{De los socios}
%=======================================================

\articulo 7. Tipos de socios.\par\nobreak

1. La asociaci\'on est\'a integrada por:

(a) Socios numerarios.

(b) Socios colaboradores.

(c) Socios honor\'{\i}ficos.

2. Son socios numerarios las personas que participan en la
asociaci\'on, bien de forma activa en los grupos de trabajo o en la
Junta Directiva, bien ofreciendo su apoyo sin participar regularmente
en los trabajos.

3. Son socios colaboradores los que de forma voluntaria realizan una
aportaci\'on econ\'omica anual por un valor de al menos el doble de la
cuota anual que le corresponda.

4. Son socios honor\'{\i}ficos las personas que hayan apoyado de forma
especialmente notoria los fines de la asociaci\'on. Ser\'an nombrados por la
Asamblea General a propuesta de la Junta Directiva o un tercio de los
socios.

5. Las personas jur\'{\i}dicas que tengan la consideraci\'on de socios
numerarios, colaboradores u honor\'{\i}ficos s\'olo podr\'an participar por
medio de un representante a lo sumo en cada uno de los grupos de trabajo y
\'organos directivos.

%------------------

\articulo 8. Adquisici\'on de la condici\'on de asociado.\par\nobreak

1. Puede ser socio cualquier persona, f\'{\i}sica o jur\'{\i}dica, que
comparta los intereses de la asociaci\'on y acepte apoyar sus actividades.
Los menores de edad deben certificar la autorizaci\'on de sus padres o
tutores legales.

2. La solicitud debe hacerse por escrito y la Junta Directiva decidir\'a
sobre su aceptaci\'on.

3. La sola solicitud implica la aceptaci\'on de los estatutos.

%------------------

\articulo 9. Derechos y deberes de los socios.\par\nobreak

1. Los socios se comprometen a apoyar los objetivos e intereses de la
asociaci\'on con los medios disponibles, adem\'as de cumplir las decisiones
y acuerdos de los \'organos de la asociaci\'on.

2. Todos los socios son iguales en deberes y derechos, con las
salvedades reguladas en este reglamento.

%------------------

\articulo 10. Cuotas de los socios.\par\nobreak

1. Las cuotas de los socios deben ser pagadas por adelantado y valen para
un a\~no natural (del uno de enero al treinta y uno de diciembre).

2. La Asamblea General decidir\'a sobre la cuant\'{\i}a de las cuotas, que
podr\'a diferenciar entre personas f\'{\i}sicas y jur\'{\i}dicas.

3.  Los socios que no hayan pagado la cuota de la asociaci\'on al terminar
el a\~no ser\'an amonestados. Tras requerimientos infructuosos, pueden
ser suspendidos de sus derechos por decisi\'on de la Junta Directiva.

4.  Los socios que, libres de deudas, se encuentren en situaci\'on
econ\'omica complicada pueden solicitar una moratoria en el pago. La
Junta Directiva decidir\'a sobre su concesi\'on y podr\'a incluso acordar
la condonaci\'on parcial o completa del pago.

5. Los socios honor\'{\i}ficos est\'an exentos del pago de las cuotas.

6. Las cuotas ya abonadas no ser\'an reembolsadas en caso de la
p\'erdida de la condici\'on de socio.

%------------------

\articulo 11. Perdida de la condici\'on de socio.\par\nobreak

La condici\'on de socio termina por:

(a) Fallecimiento de personas f\'{\i}sicas o disoluci\'on de personas
jur\'{\i}dicas.

(b) Baja voluntaria, que puede realizarse en cualquier momento mediante
escrito a la Junta Directiva.

(c) Falta de pago de dos cuotas anuales.

(d) Expulsi\'on, acordada por la Junta Directiva, por causa grave y pese a
amonestaci\'on previa. El socio expulsado puede interponer recurso ante la
siguiente Asamblea General. Hasta la plena efectividad de la expulsi\'on
quedan suspendidos los derechos del socio expulsado.

\bigskip

%=======================================================
\titulo{III}{De los \'organos de gobierno}
%=======================================================

\articulo 12. Los \'organos de gobierno de Cervan\TeX{} son:\par\nobreak

1. La Asamblea General.

2. La Junta Directiva.

3. El Presidente.

4. El Tesorero.

5. El Secretario.

%------------------

\capitulo{1}{La Asamblea General}

%------------------

\articulo 13. La Asamblea General estar\'a constituida por todas las personas
f\'{\i}sicas y un representante de cada persona jur\'{\i}dica miembros de
Cervan\TeX, y ser\'a su m\'aximo \'organo rector.

%------------------

\articulo 14. Son competencias de la Asamblea General:\par\nobreak

1. Aprobar los estatutos de Cervan\TeX{} y sus posibles modificaciones.

2. Establecer las l\'{\i}neas generales de actuaci\'on de Cervan\TeX{}.

3. Elegir al Presidente de la Junta Directiva.

4. Ratificar al Secretario y al Tesorero, previa propuesta del Presidente.

5. Determinar el n\'umero de vocales de la Junta Directiva.

6. Aprobar o rechazar las mociones de censura.

7. Aprobar el presupuesto anual de Cervan\TeX.

8. Conocer anualmente y aprobar, en su caso, los balances y estados de
cuentas de Cervan\TeX.

9. Aprobar anualmente, a propuesta de la Junta Directiva, las cuotas que
deber\'an abonar los socios numerarios.

10. Determinar las condiciones en las que los socios disfrutar\'an de los
beneficios derivados de su condici\'on de tales.

11. Aprobar la disoluci\'on de Cervan\TeX.

12. Resolver cuantos asuntos sean sometidos a su consideraci\'on por la
Junta Directiva.

13. Aprobar las normativas que permitan la estandarizaci\'on del uso de
\TeX{} en espa\~nol.

%------------------

\articulo 15. La Asamblea General se reunir\'a con car\'acter ordinario una vez
al a\~no, y con car\'acter extraordinario, cuando as\'{\i} lo acuerde la Junta
Directiva o lo solicite el veinte por ciento de los miembros de Cervan\TeX{}.
En este \'ultimo caso, la petici\'on, que incluir\'a la cuesti\'on o
cuestiones a tratar, deber\'a dirigirse por escrito al Presidente, quien
ordenar\'a la convocatoria para una fecha en los treinta d\'{\i}as siguientes.
En todo caso tendr\'an car\'acter extraordinario las convocatorias para tratar
alguno de los puntos siguientes:

1. Modificaci\'on de los Estatutos.

2. Disoluci\'on de Cervan\TeX.

%------------------

\articulo 16. Las sesiones de la Asamblea General, tanto ordinarias como
extraordinarias, ser\'an convocadas por carta o procedimiento electr\'onico
un mes antes de la fecha de su celebraci\'on. La convocatoria incluir\'a lugar,
d\'{\i}a y hora de la reuni\'on y el orden del d\'{\i}a con los asuntos que se
tratar\'an.

%------------------

\articulo 17. Las sesiones de la Asamblea General ser\'an presididas por el
Presidente de Cervan\TeX{} y el Secretario el actuar\'a como tal.

Sus acuerdos se tomar\'an por mayor\'{\i}a simple, salvo que legal o
estatuariamente se exija mayor proporci\'on, admiti\'endose la delegaci\'on
de voto y el voto por correo (tanto ordinario como electr\'onico).
Cada socio podr\'a llevar a la Asamblea General un m\'aximo de~5 votos
delegados, con la salvedad que se indica en el art\'{\i}culo siguiente.

La Asamblea General Ordinaria quedar\'a v\'alidamente constituida con la
presencia de la mitad m\'as uno de sus miembros, en primera convocatoria, y
cualquiera que sea el n\'umero de los presentes, en segunda convocatoria,
media hora m\'as tarde.

%------------------

\articulo 18. Las reuniones de la Asamblea General no deber\'an ser,
necesariamente, presenciales. Podr\'an celebrarse a distancia con la
ayuda de los adecuados procedimientos electr\'onicos de comunicaci\'on.
% (L\'ease multiconferencia, irc. \dots)

As\'{\i} mismo, las votaciones podr\'an celebrarse a distancia,
electr\'onicamente. En este caso, no se admitir\'an votos delegados
transmitidos electr\'onicamente.

En todo caso, se habilitar\'an mecanismos de control para verificar
la autenticidad de las personas participantes en la asamblea y
de sus votaciones.
% (Cosas del estilo de pgp, si es que realmente se ve necesidad
% de hacer algo de esto)

%------------------

\capitulo{2}{La Junta Directiva}

%------------------

\articulo 19. La Junta directiva estar\'a compuesta por el Presidente y,
al menos, dos vocales, actuando uno de ellos como Secretario y otro como
Tesorero. El n\'umero de vocales ser\'a determinado por la Asamblea General.

\'Estos ser\'an elegidos como se indica en los art\'{\i}culos 23 y
siguientes.

%------------------

\articulo 20. Son competencias de la Junta Directiva:\par\nobreak

1. Aprobar las actividades a desarrollar por Cervan\TeX{} en el marco de
las l\'{\i}neas de actuaci\'on establecidas por la Asamblea General.

2. Adoptar acuerdos sobre aquellos asuntos que no est\'en expresamente
atribuidos a la Asamblea General por estos Estatutos.

3. Elaborar el proyecto de presupuestos de Cervan\TeX{} y someterlo a la
Asamblea General.

4. Elaborar la propuesta de cuotas anuales.

5. Elaborar la memoria anual de actividades, as\'{\i} como el balance y el
estado de cuentas y someterlas a la Asamblea General para su aprobaci\'on.

6. Estudiar, y aprobar en su caso, convenios de colaboraci\'on
e intercambio con otras sociedades.

7. Aprobar la creaci\'on grupos de trabajo y coordinar el trabajo de
estos entre s\'{\i}.

8. Refrendar, a propuesta del Presidente, el nombramientos del
Secretario y Tesorero de Cervan\TeX{}.

9. Aprobar, previa petici\'on, la moratoria o condonaci\'on total
o parcial de la cuota de alg\'un socio.

%------------------

\articulo 21. La Junta Directiva se reunir\'a al menos dos veces al a\~no y
siempre que lo solicite alguno de sus miembros, en escrito dirigido al
Presidente. En este \'ultimo caso, la fecha de la reuni\'on ser\'a uno de los
quince d\'{\i}as siguientes.

Las sesiones de la Junta Directiva ser\'an convocadas por el Secretario por
orden del Presidente.

La Junta Directiva quedar\'a v\'alidamente constituida con la presencia de,
al menos, la mitad m\'as uno de sus miembros, tanto en primera como
en segunda convocatoria, media hora m\'as tarde. Sus acuerdos se tomar\'an
por mayor\'{\i}a simple de votos. En caso de empate,
decidir\'a el voto del Presidente.

%------------------

\articulo 22. Los miembros de la Junta Directiva ser\'an elegidos por tres
a\~nos, pudiendo presentarse a la reelecci\'on. No obstante, ning\'un miembro
de Cervan\TeX{} podr\'a formar parte de la Junta Directiva m\'as
de seis a\~nos consecutivos, con la salvedad de la excepci\'on mencionada
en el art\'{\i}culo~24.

%------------------

\articulo 23. La renovaci\'on de la Junta Directiva se llevar\'a a cabo
en una sesi\'on de la Asamblea General.

Una vez hecha p\'ublica la convocatoria para dicha sesi\'on, la
Junta Directiva abrir\'a un plazo para la presentaci\'on de candidatos;
este plazo ser\'a, como m\'aximo, de tres meses.

Habr\'a candidaturas separadas para Presidente y para vocales de la
Junta Directiva. Los miembros que deseen presentar su candidatura deber\'an
dirigirse con tal fin, por escrito, al Presidente.

Transcurrido el plazo se\~nalado, la Junta Directiva aprobar\'a la lista de
candidatos que ser\'a remitida a todos los socios con la suficiente
antelaci\'on para que puedan, en su caso, enviar el voto por correo.

%------------------

\articulo 24. Si, en el plazo establecido, no se presentan suficientes
candidatos, se abrir\'a un nuevo plazo de igual duraci\'on y con igual
procedimiento de elecci\'on.
En este nuevo plazo, podr\'an presentarse candidatos que ya hayan permanecido
seis a\~nos consecutivos en la Junta Directiva.

%------------------

\articulo 25. Una vez elegido el Presidente, \'este propondr\'a al
Secretario y al Tesorero de entre los vocales electos, que deber\'an ser
refrendados por la Asamblea General. En el caso de que la Asamblea
General rechazase la propuesta, el Presidente propondr\'a, si los
hubiere, otros vocales para estos cargos. Si todas las propuestas son
rechazadas, se escoger\'a Secretario y Tesorero por elecci\'on directa
de la Asamblea General, actuando \'estos como vocales en la nueva Junta
Directiva.

\articulo 26. Si se produce la dimisi\'on del Presidente, o su cese por
muerte o enfermedad grave, la Junta Directiva continuar\'a ejerciendo
interinamente sus funciones hasta que se produzca una nueva
elecci\'on de Presidente y Junta Directiva por el procedimiento ordinario
descrito en los art\'{\i}culos anteriores.

\articulo 27. Si se produce el cese de alguno de los vocales, el Presidente
convocar\'a, en el plazo de tres meses, elecciones parciales para la plaza
vacante.

Si el vocal que ha cesado act\'ua como Tesorero o Secretario, se proceder\'a
a la renovaci\'on del cargo como indica el art\'{\i}culo~25.

\articulo 28. Mociones de censura:\par\nobreak

Se podr\'an presentar mociones de censura a la Junta Directiva.
La moci\'on se presentar\'a ante el Presidente y deber\'a incluir
una propuesta de nuevo Presidente y de tantos vocales como
tenga la Junta Directiva.

La moci\'on de censura se debatir\'a, en un plazo m\'{\i}nimo de quince
d\'{\i}as y m\'aximo de tres meses, en asamblea ordinaria o extraordinaria.
Si, en el citado plazo, no hay ninguna asamblea ordinaria prevista,
el Presidente deber\'a convocar asamblea extraordinaria con ese punto
en el Orden del D\'{\i}a. 

La moci\'on de censura se considerar\'a aprobada si obtiene mayor\'{\i}a
simple de la Asamblea General.

S\'olo se podr\'a presentar una moci\'on de censura por a\~no.

No se podr\'an presentar mociones de censura cuando haya elecciones
convocadas.

%------------------

\capitulo{3}{El Presidente}

%------------------

\articulo 29. El Presidente ser\'a elegido seg\'un se indica en los
art\'{\i}culos 23 y siguientes de estos Estatutos.

%------------------

\articulo 30. Son funciones del Presidente:\par\nobreak

1. Ostentar la representaci\'on legal y oficial de Cervan\TeX{} ante
Autoridades, Corporaciones y Organismos, tanto p\'ublicos como privados, y
personas de toda clase y naturaleza y, con tal car\'acter, firmar escrituras
publicas, contratos privados y dem\'as documentos relacionados con
Cervan\TeX{}, previo acuerdo de la Junta Directiva.

2. Convocar y presidir las reuniones de la Junta Directiva y de la
Asamblea General, tanto ordinarias como extraordinarias, y poner el Visto
Bueno a las actas correspondientes.

3. Vigilar el cumplimiento de los acuerdos adoptados por la Junta de
Directiva y la Asamblea General.

4. Ordenar y dirigir los debates en las reuniones de la Junta Directiva y
de la Asamblea General y decidir los empates en las votaciones de la Junta
Directiva.

5. Proponer a la Asamblea General el Secretario y Tesorero de Cervan\TeX{}.

6. Promover ante la Asamblea General la renovaci\'on parcial de la Junta
Directiva si se produce la dimisi\'on o cese (por muerte o enfermedad grave)
de alguno de sus miembros.

7. Realizar, de acuerdo con la Junta Directiva, las gestiones que
considere convenientes para el funcionamiento de Cervan\TeX{} y
el cumplimiento de sus fines.

8. Comunicarse con el Editor General para conocer el estado de
las publicaciones de Cervan\TeX{} y hacer las propuestas que considere
pertinentes.

9. La coordinaci\'on, en primera instancia, de los grupos de trabajo.

10. El presidente podr\'a delegar, previa aprobaci\'on de la Junta Directiva,
alguna de estas funciones en los vocales o en otros miembros de Cervan\TeX{}.

%------------------

\capitulo{4}{El Tesorero}

%------------------

\articulo 31. El Tesorero ser\'a nombrado por la Asamblea General
de entre los vocales de la Junta Directiva, a propuesta del Presidente,
seg\'un se indica en el art\'{\i}culo 23 de estos Estatutos.

%------------------

\articulo 32. Son funciones del Tesorero las siguientes:\par\nobreak

1. La recaudaci\'on de las cuotas de los socios.

2. La recaudaci\'on de todos los ingresos que Cervan\TeX{} pueda tener.

3. Efectuar los pagos acordados por la Junta Directiva.

4. Presentar el balance econ\'omico anual a la Asamblea General.

%------------------

\capitulo{5}{El Secretario}

%------------------

\articulo 33. El Secretario ser\'a nombrado por la Asamblea General
de entre los vocales de la Junta Directiva, a propuesta del Presidente,
seg\'un se indica en el art\'{\i}culo 23 de estos Estatutos.

%------------------

\articulo 34. Son funciones del Secretario las siguientes:\par\nobreak

1. Asumir las funciones del Presidente en casos de ausencia o enfermedad.

2. Actuar como secretario en todas las reuniones de la Junta Directiva y
de la Asamblea General, levantando acta de las mismas. Para ello habilitar\'a
el correspondiente libro de actas y se encargar\'a de su custodia.

3. Archivar y custodiar los documentos de Cervan\TeX.

4. Elaborar y actualizar el registro de miembros de Cervan\TeX.

5. Registrar y cursar la correspondencia de Cervan\TeX.

6. Extender las certificaciones que correspondan, con el Visto Bueno del
Presidente.

7. Tener actualizada un acta en la que se recojan todas las actividades
oficiales de Cervan\TeX.

8. Informar a los miembros de Cervan\TeX{} sobre aquellas cuestiones
relativas al funcionamiento de la misma en las que est\'en interesados.

9. Dirigir al personal auxiliar que trabaje para Cervan\TeX.

\bigskip

%=======================================================
\titulo{IV}{Del Editor General}
%=======================================================

\articulo 35. El Editor General ser\'a nombrado por los miembros del
grupo de trabajo encargado de las publicaciones.

%------------------

\articulo 36. Son funciones del Editor General las siguientes:\par\nobreak

1. Se responsabilizar\'a de la edici\'on y distribuci\'on de las revistas y
publicaciones de Cervan\TeX{} con ayuda de su grupo de trabajo.

2. Procurar\'a formalizar acuerdos e intercambios con las publicaciones de
otras sociedades afines y de editoriales cient\'{\i}ficas.

3. Mantener informado al Presidente de Cervan\TeX{} del trabajo realizado.

\bigskip

%=======================================================
\titulo{V}{Del r\'egimen econ\'omico}
%=======================================================

\articulo 37. Cervan\TeX{} se constituye sin \'animo de lucro y sus
recursos econ\'omicos no exceder\'an del tope legal establecido.
El l\'{\i}mite del presupuesto anual de Cervan\TeX{} ser\'a
de~65536 euros.

%------------------

\articulo 38. Cervan\TeX{} no cuenta con patrimonio fundacional.

%------------------

\articulo 39. Los medios de que dispone Cervan\TeX{} para el
cumplimiento de sus fines son los siguientes:

(a) Cuotas de sus socios.

(b) Subvenciones que se obtengan.

(c) Donaciones que se reciban.

(d) Renta y rendimiento de su patrimonio.

(e) Ingresos que produzcan sus actividades.

\bigskip

%=======================================================
\titulo{VI}{Modificaci\'on de los Estatutos}
%=======================================================

\articulo 40. Para introducir alguna modificaci\'on en los Estatutos de
Cervan\TeX{} har\'a falta una propuesta firmada por al menos el veinte
por ciento de sus miembros, o por la Junta Directiva, que deber\'a ser
enviada al Presidente. Este deber\'a convocar Asamblea General
Extraordinaria para dicho prop\'osito. Para la aprobaci\'on de una
enmienda se requerir\'a el voto favorable de dos tercios de los votantes.

\bigskip

%=======================================================
\titulo{VII}{Disoluci\'on de Cervan\TeX}
%=======================================================

\articulo 41. Para disolver Cervan\TeX{} har\'a falta una propuesta firmada
por al menos el veinte por ciento de sus miembros, o por la Junta
Directiva, que deber\'a ser enviada al Presidente. Este deber\'a convocar
Asamblea General Extraordinaria para dicho prop\'osito. Para la
disoluci\'on de Cervan\TeX{} se requerir\'a el voto favorable de
dos tercios de los votantes. El patrimonio resultante, si lo hubiere,
se donar\'a a una Sociedad o Instituci\'on con fines an\'alogos.

\bye


